% Figure: schematic machinery vibration spectrum, log amplitude versus
% frequency in orders of running speed, with labelled fault peaks.
% Synthetic spectrum built from a few Gaussian-like peaks on a noise floor.
\begin{figure}[htbp]
\centering
\begin{tikzpicture}
\begin{axis}[
  width=12cm, height=6.5cm,
  xlabel={frequency (orders of running speed, $1\times = f_{\mathrm{run}}$)},
  ylabel={amplitude (dB, ref.)},
  xmin=0, xmax=6, ymin=-5, ymax=42,
  axis lines=left,
  domain=0.05:6, samples=600,
  every axis x label/.style={at={(current axis.right of origin)}, anchor=west},
  every axis y label/.style={at={(current axis.above origin)}, anchor=south},
  grid=major, grid style={enggray!20},
  tick label style={font=\scriptsize},
  label style={font=\small},
]
% noise floor + peaks. Peaks: 1x (imbalance, 38), 2x (misalign, 30),
% 0.5x (looseness subharm, 18), broad band bump near 4.3x (bearing, 22).
\addplot[engblue, thick] {
  3
  + 35*exp(-((x-1.0)^2)/(2*0.015))
  + 27*exp(-((x-2.0)^2)/(2*0.02))
  + 15*exp(-((x-0.5)^2)/(2*0.012))
  + 19*exp(-((x-4.3)^2)/(2*0.25))
  + 12*exp(-((x-3.0)^2)/(2*0.02))
};
\node[engred, font=\scriptsize, anchor=south] at (axis cs:1.0,39) {$1\times$ imbalance};
\node[engamber, font=\scriptsize, anchor=south] at (axis cs:2.0,31) {$2\times$ misalign};
\node[engorange, font=\scriptsize, anchor=south] at (axis cs:0.5,19) {$0.5\times$ loose};
\node[enggray, font=\scriptsize, anchor=south] at (axis cs:4.3,23) {bearing band};
\end{axis}
\end{tikzpicture}
\caption[A machinery vibration signature]{A schematic vibration
spectrum of a rotating machine, amplitude in decibels against
frequency expressed in orders of the running speed ($1\times$ is one
rotation per cycle). The discipline of signature analysis reads the
spectrum as a fault diagnosis. A dominant peak at $1\times$ marks
rotor imbalance; a peak at $2\times$ comparable to or larger than
$1\times$ marks shaft misalignment; a subharmonic peak at $0.5\times$
marks mechanical looseness; a broad band of energy at
non-integer-order frequencies (here near $4.3\times$) marks rolling
element bearing defects, whose frequencies follow from the bearing
geometry and are rarely integer multiples of the running speed. A
resonance shows as a sharp peak that stays fixed in absolute
frequency as the running speed changes, so it migrates across the
order axis. Trending these amplitudes over months is the basis of
condition-based maintenance.}
\label{fig:vol03ch06:spectrum}
\end{figure}
